\documentclass{article}
\usepackage{nicematrix}
\usepackage{siunitx}
\usepackage{makecell}
\usepackage{multirow,bigdelim}
\usepackage{tikz}
\usetikzlibrary{decorations.pathreplacing,calc,tikzmark}

\begin{document}

\begin{center}
\begin{tabular}{c|c|c||c|c|c}
\multicolumn{6}{c}{TABLE VII}\\[3pt]
\hline
\hline
\rule{0pt}{1.5\normalbaselineskip}
$n$ & $4.917\!\times\!{n}$ & \makecell{\small{Observed}\\\small{Charge}} & $n$ & $4.917\!\times\!{n}$ & \makecell{\small{Observed}\\\small{Charge}}\\[8pt]
\hline
\rule{0pt}{1\normalbaselineskip}
1 & 4.917 & $\cdots$ & 10 & 49.17 & 49.41\\
2 & 9.834 & $\cdots$ & 11 & 54.09 & 53.91\\
3 & 14.75 & $\cdots$ & 12 & 59.00 & 59.12\\
4 & 19.66 & 19.66  & 13 & 63.92 & 63.68\\
5 & 24.59 & 24.60  & 14 & 68.84 & 68.65\\
6 & 29.50 & 29.62  & 15 & 73.75 & $\cdots$\\
7 & 34.42 & 34.47  & 16 & 78.67 & 78.34\\
8 & 39.34 & 39.38  & 17 & 83.59 & 83.22\\
9 & 44.25 & 44.42  & 18 & 88.51 & $\cdots$\\[3pt]
\hline
\end{tabular}
\end{center}

\bigskip

\begin{center}
\small
\begin{tabular}{|S|c||S|c|c|c|c|c|c|c|}
\hline
\rule{0pt}{1.5\normalbaselineskip}
 & $d$ = 0.5cm & & $d$ = 0.5cm & \makecell{Charge\\on ion} & & & \makecell{Frictional\\charge} & & \\[10pt]
\hline
\rule{0pt}{1.5\normalbaselineskip}
$t_g$ & \makecell{$v_1(=d/t_g)$\\(cm/sec)} & $t_F$ & \makecell{$v_2(=d/t_F)$\\(cm/sec)} & $(v_{2}'\!-\!v_2)$ & $n'$ & $\frac{v_{2}'-v_2}{n'}$ & $v_1+v_2$ & $n$ & $\frac{v_1+v_2}{n}$\\[10pt]
\hline
\rule{0pt}{1\normalbaselineskip}
18.2    & .00286 & 3.8 & .01316 & & & & .01602 & 3 & .00534\\
18.6    & \emph{avr}  & & & .00470 & 1 & .00470 & & & \\
19.2    & & 2.8 & .01786 & & & & & & \\
18.0    & & & & .01561 & 3 & .00520 & & & \\
17.2    & & 22.2 & .00225 & & & & & & \\
15.4    & & & & .00544 & 1 & .00544 & & & \\
16.7    & & 6.5 & .00769 & & & & & & \\
18.0    & & & & .00541 & 1 & .00541 & & & \\
15.4    & & 21.9 & .00228 & & & & & & \\
17.3    & & & & .01123 & 2 & .00562 & & & \\
18.4    & & 3.7 & .01351 & & & & & & \\[3pt]
\hline
\rule{0pt}{1\normalbaselineskip}
17.5 & & & & & & .00527 & & & .00534 \\
\emph{avr} & \makecell{test} & & & & & \emph{avr} & & &\\[3pt]
\hline
\end{tabular}
\end{center}

\bigskip

\begin{table}[htp]
\centering
\begin{minipage}{\textwidth}
\centering
\begin{tabular}{c|r@{\hspace{10pt}}|l@{\hspace{10pt}}|c|c|c||c|c|c}
\multicolumn{9}{c}{TABLE VI\footnote{[The bracketed numbers are our corrections of errors in the original paper.]}}\\[3pt]
\hline
\makecell{$t_g$\\[-2pt]\footnotesize{Sec.}} & \makecell{$t_{\scriptstyle{F}}$\\[-2pt]\footnotesize{Sec.}} & $\ \ \frac{1}{t_F}$ & $\frac{1}{t'_F}-\frac{1}{t_F}$ & $n'$ & $\frac{1}{n'}(\frac{1}{t'_F}-\frac{1}{t_F})$ & $\frac{1}{t_g}+\frac{1}{t_F}$ & n & $\frac{1}{n}(\frac{1}{t_g}+\frac{1}{t_F})$\\[5pt]
\hline
\rule{0pt}{1\normalbaselineskip}%space before first row numbers
11.848 & 80.708 & .01236\phantom{0}\tikzmark{3} & & & & .09655 & 18 & .005366\\
11.890 & 22.366\tikzmark{1} & & .03234 & 6 & .005390 & & & \\
11.908 & 22.390 & .04470\phantom{0}\tikzmark{4}\tikzmark{6} & & & & .12887 & 24 & .005371\\
11.904 & 22.368\tikzmark{2} & & .03751 & 7 & .005358 & & & \\
11.882 & 140.565 & .007192\tikzmark{5}\tikzmark{7}  & & & & .09138 & 17 & .005375\\
& & & .005348 & 1 & .005348 & & &\\
11.906 & 79.600 & .01254\phantom{0}\tikzmark{8} & & & & .09673 & 18 & .005374\\
11.838 & 34.748\tikzmark{9} & & .01616  & 3 & .005387 & & & \\
11.816 & 34.762 & .02870  & & & & .11289 & 21 & .005376\\
11.776 & 34.846\tikzmark{10} & & & & & & \\
\hline
\end{tabular}
\begin{tabular}{l@{\hskip 3em} l l@{\hskip 5em} l l}\\[-8pt]
& Duration of exp. & = 45 min. & Pressure & = 75.62 cm.\\
& Plate distance & = 16 mm. & Oil density & = .9199\\
& Fall distance & = 10.21 mm. & Air viscosity & = $1,824\times 10^{-7}$ [poise]\\
\end{tabular}
\begin{tikzpicture}[remember picture]
% Column 2 braces
\draw[overlay,decorate,
    decoration = {brace,raise=3pt,amplitude=4pt,aspect=0.45}]([shift={(0pt,2pt)}]pic cs:1) -- ([shift={(0pt,2pt)}]pic cs:2);
\draw[overlay,decorate,
    decoration = {brace,raise=3pt,amplitude=4pt,aspect=0.45}]([shift={(0pt,3pt)}]pic cs:9) -- ([shift={(0pt,2pt)}]pic cs:10);
% Column 3 braces
\draw[overlay,decorate,
    decoration = {brace,raise=2pt,amplitude=4pt,aspect=0.45}]([shift={(0pt,3pt)}]pic cs:3) -- ([shift={(0pt,3pt)}]pic cs:4);
\draw[overlay,decorate,
    decoration = {brace,raise=2pt,amplitude=4pt,aspect=0.45}]([shift={(0pt,3pt)}]pic cs:6) -- ([shift={(0pt,3pt)}]pic cs:5);
\draw[overlay,decorate,
    decoration = {brace,raise=2pt,amplitude=4pt,aspect=0.45}]([shift={(0pt,3pt)}]pic cs:7) -- ([shift={(0pt,2pt)}]pic cs:8);

\end{tikzpicture}
\end{minipage}
\end{table}

\begin{tabular}{r r r r}
\hline
1 & 2\tikzmark{A} & 3 & 4\\
5 & 6\tikzmark{D} & \tikzmark{B}7 & 8\\
9 & 10\tikzmark{C} & 11 & 12\\
\hline
\end{tabular}
\tikz[remember picture]
\draw[overlay,decorate,
    decoration = {brace,raise=6pt,amplitude=3pt,aspect=0.45}]([shift={(0pt,2pt)}]pic cs:A) -- ([shift={(0pt,2pt)}]pic cs:C);
\bigskip

\begin{tabular}{r r r r}
\hline
1 & 2\tikzmark{a} & 3 & 4\\
5 & 6\tikzmark{d} & \tikzmark{b}7 & 8\\
9 & 10\tikzmark{c} & 11 & 12\\
\hline
\end{tabular}
\tikz[remember picture]\draw[overlay]([shift={(1pt,2pt)}]pic cs:a) -- ([shift={(-2pt,4pt)}]pic cs:b);
\tikz[remember picture]\draw[overlay]([shift={(1pt,2pt)}]pic cs:c) -- ([shift={(-2pt,4pt)}]pic cs:b);

\bigskip

%\url{https://tex.stackexchange.com/a/1570/86}

\hfill\tikzmark{right}
\begin{itemize}
\item First line
\item Second line \tikzmark{2nd}
\item Third line, which is quite long and seemingly tedious in the extreme
\item Fourth line, which isn't as long as the third \tikzmark{4th}
\item Fifth line
\end{itemize}

\begin{tikzpicture}[overlay, remember picture]
\node[anchor=base] (a) at (pic cs:2nd) {\vphantom{h}}; % push the mark to the top of the line (ie including ascenders)
\node[anchor=base] (b) at (pic cs:4th) {\vphantom{g}}; % push the mark to the bottom of the line (ie including descenders)
\draw [decoration={brace,amplitude=0.5em},decorate,ultra thick,gray]
 (a.north -| {pic cs:right}) -- (b.south -| {pic cs:right});
\end{tikzpicture}

\begin{equation}
P(\pmb{a}, \pmb{b}) = \int d\lambda\, \rho(\lambda) A(\pmb{a}, \lambda) B(\pmb{b}, \lambda).\footnote{[According to the hypothesis Bell is investigating, the "more complete specification" of the particles' state, the extra information that would tell the observer who knew it which way the particle would respond to a given measurement, is given by $\lambda$. Equation (2) considers all the possible $\lambda$'s and weights them according to how likely each is. Mathematically, that means integrating over all the possible values of $\lambda$ and weighting each by its likelihood, represented by the probability distribution $\rho(\lambda)$. It may be helpful to note that $P$ was likely chosen to stand for "product," since the term we are integrating $\rho$ against is the product of $A$ and $B$. Because we've defined $A$ and $B$ such that the results of measurements yield the values $+1$ and $-1$, the product for any given value of $\lambda$ will always be either $+1$ (meaning that the two particles would both be measured to be pointing either in the selected directions $\pmb{a}$ and $\pmb{b}$ or in directions opposed to them) or $-1$ (meaning that the directions of $\pmb{\sigma}_1$ and $\pmb{\sigma}_2$ would be measured to be opposed, whether because $\pmb{\sigma}_1$ would \emph{not} be measured to point in direction $\pmb{a}$ and $\pmb{\sigma}_2$ \emph{would} be measured to point in direction $\pmb{b}$ or because $\pmb{\sigma}_1$ \emph{would} be measured to point in direction $\pmb{a}$ and $\pmb{\sigma}_2$ would \emph{not} be measured to point in direction $\pmb{b}$. Since it is weighted sum of all the positive and negative terms, $P(\pmb{a},\pmb{b})$ is thus a measure of the \emph{correlation} between the spin-direction of the two particles relative to the angles $\pmb{a}$ and $\pmb{b}$.]}
\end{equation}



\end{document}
